\chapter{Concussion (Mild Traumatic Brain Injury)}

\section*{Definition}
A concussion is a mild traumatic brain injury caused by a blow, bump, or jolt to the head or body that transmits force to the brain. Symptoms may begin immediately or hours later. Most recover, but urgent assessment is needed for danger signs or worsening symptoms.

\section*{When to See a Doctor}
\begin{itemize}
\item Headache or pressure in the head
\item Dizziness or balance problems
\item Nausea or vomiting
\item Blurred or double vision
\item Sensitivity to light or noise
\item Confusion, disorientation, or feeling foggy
\item Difficulty concentrating or remembering
\item Slowed reaction times
\item Mood changes (irritability, sadness, anxiety)
\item Sleep disturbances (insomnia, drowsiness)
\item Loss of consciousness (occurs in less than 10 percent of concussions)
\end{itemize}

\section*{How It Develops}
The sudden acceleration-deceleration force causes shearing of axons, disruption of neuronal membranes, and a complex neurometabolic cascade. There is an influx of calcium and efflux of potassium, impaired oxidative metabolism, decreased cerebral blood flow, and altered neurotransmitter release. These changes produce the clinical symptoms of concussion, which are functional rather than structural.

\section*{Causes}
\begin{itemize}
\item Sports injuries (football, soccer, hockey, boxing, rugby)
\item Falls (most common cause in children and elderly)
\item Motor vehicle accidents
\item Assault or physical violence
\item Workplace injuries
\item Blast injuries in military settings
\end{itemize}

\section*{Related Symptoms}
\begin{itemize}
\item Head injury
\item Headache (post-concussive)
\item Dizziness
\item Memory loss
\item Neck pain (whiplash-associated)
\item Post-concussion syndrome
\end{itemize}

\section*{Medical Evaluation}
\begin{itemize}
 \item Clinical evaluation, with age-appropriate symptom and neurologic assessment tools when available
\item Neurological examination
 \item CT scan when decision rules and clinical findings indicate possible intracranial injury; concussion alone is not an automatic indication for imaging
\item MRI for persistent symptoms
\item Neuropsychological testing for return-to-play decisions
\end{itemize}

\section*{Treatment Options}
\begin{itemize}
 \item Relative physical and cognitive rest for no more than 1--2 days, followed by gradual return to activities that do not substantially worsen symptoms
\item Gradual return to activity protocol (school, work, sport)
\item Symptom management (headache, sleep, mood)
\item Referral to concussion specialist for persistent symptoms
\item Vestibular therapy for balance issues
\item Avoid repeat concussion risk during recovery
\end{itemize}
\section*{Self-Care and Home Management}
\begin{itemize}
 \item Reduce activity for the first 1--2 days, then gradually resume light activity as tolerated
 \item Limit screens or other activities only when they worsen symptoms; complete avoidance is not routinely required
 \item Avoid contact sports and activities with risk of another head injury until medically cleared
\item Get adequate sleep
\item Avoid alcohol and recreational drugs
 \item Use pain medicine only as directed and according to clinician advice, especially if bleeding risk or another injury is possible
\item Return to activities gradually under medical supervision
\end{itemize}

\section*{References}
\begin{thebibliography}{99}
\bibitem{concussion:ref1} McCrory P, Meeuwisse W, Dvorak J, et al. ``Consensus Statement on Concussion in Sport: The 5th International Conference on Concussion in Sport.'' Br J Sports Med. 2017;51(11):838--847.
\bibitem{concussion:ref2} Giza CC, Hovda DA. ``The Neurometabolic Cascade of Concussion.'' J Athl Train. 2001;36(3):228--235.
\bibitem{concussion:ref3} Lumba-Brown A, Yeates KO, Sarmiento K, et al. ``CDC Guideline on the Diagnosis and Management of Mild Traumatic Brain Injury Among Children.'' JAMA Pediatr. 2018;172(11):e182853.
\bibitem{concussion:ref4} Leddy JJ, Haider MN, Ellis MJ, et al. ``Practice Guideline: Treatment of Concussion and Mild Traumatic Brain Injury.'' Neurology. 2023;100(16):770--779.
\bibitem{concussion:ref5} Broglio SP, Cantu RC, Gioia GA, et al. ``National Athletic Trainers' Association Position Statement: Management of Sport Concussion.'' J Athl Train. 2014;49(2):245--265.
\bibitem{concussion:ref6} Meehan WP III, Bachur RG. ``Sport-Related Concussion.'' Pediatrics. 2009;123(1):114--123.
\end{thebibliography}
